\section*{Hoofdstuk \ref{ch:g8:balanced-equations} --- Behoud van massa en kloppende reactievergelijkingen}

\begin{solution}{exo:g8:balanced-equations:1}
Bij een \omterm{def:g7:chemical-reactions:reaction}{chemische reactie} is de totale massa van de gevormde
\omterm{def:g7:chemical-reactions:reactant}{reactieproducten} gelijk aan de totale massa van de verbruikte
\omterm{def:g7:chemical-reactions:reactant}{beginstoffen}.
\end{solution}

\begin{solution}{exo:g8:balanced-equations:2}
De \omterm{def:g8:balanced-equations:coefficient}{coëfficiënt} is 3. Het stelt 3 koolstofatomen en $3 \times 2 = 6$
zuurstofatomen voor.
\end{solution}

\begin{solution}{exo:g8:balanced-equations:3}
Ja: aan elke kant 1 C en 2 O.
\end{solution}

\begin{solution}{exo:g8:balanced-equations:4}
\ce{2Mg + O2 -> 2MgO}: aan elke kant 2 Mg en 2 O.
\end{solution}

\begin{solution}{exo:g8:balanced-equations:5}
Stikstof: links 2 N, dus 2 \ce{NH3}; dan rechts 6 H, dus 3
\ce{H2}: \ce{N2 + 3H2 -> 2NH3}.
\end{solution}

\begin{solution}{exo:g8:balanced-equations:6}
Volgens de wet van \omterm{def:g8:balanced-equations:conservation}{behoud van massa}: $44 - 12 = \qty{32}{g}$ zuurstof.
\end{solution}

\begin{solution}{exo:g8:balanced-equations:7}
Het gevormde gas ontsnapt in het lokaal, dus de balans weegt het niet
meer. Er is geen massa vernietigd: de ontbrekende massa zit in de \omterm{def:g4:air-a-mixture-of-gases:air}{lucht}
van het lokaal.
\end{solution}

\begin{solution}{exo:g8:balanced-equations:8}
Koolstof: aan elke kant 1. Waterstof: links 4, dus 2 \ce{H2O}. Zuurstof:
rechts $2 + 2 = 4$, dus 2 \ce{O2}:
\ce{CH4 + 2O2 -> CO2 + 2H2O}.
\end{solution}

\begin{solution}{exo:g8:balanced-equations:9}
Als je \ce{H2O} in \ce{H2O2} verandert, verander je de stof: \ce{H2O2}
is waterstofperoxide, geen water. Alleen de \omterm{def:g8:balanced-equations:coefficient}{coëfficiënten} mogen
veranderen: \ce{2H2 + O2 -> 2H2O}.
\end{solution}

\begin{solution}{exo:g8:balanced-equations:10}
$20 \times 5 = 100$ zuurstofmoleculen; $20 \times 4 = 80$
watermoleculen.
\end{solution}

\begin{solution}{exo:g8:balanced-equations:11}
Met 2 \ce{C2H6}: 4 C, dus 4 \ce{CO2}; 12 H, dus 6 \ce{H2O}; zuurstof
rechts $8 + 6 = 14$, dus 7 \ce{O2}:
\ce{2C2H6 + 7O2 -> 4CO2 + 6H2O}.
\end{solution}

\begin{solution}{exo:g8:balanced-equations:12}
Nog steeds \qty{850.0}{g}. De doos is afgesloten: de \qty{0.4}{g}
kaarsvet die verbrand is, is koolstofdioxide en waterdamp geworden, die
nog in de doos zitten, samen met de zuurstof die ze hebben opgenomen.
Er is niets binnengekomen of weggegaan.
\end{solution}

\begin{solution}{pb:g8:balanced-equations:1}
\textbf{1.} Van de zuurstof uit de \omterm{def:g4:air-a-mixture-of-gases:air}{lucht}, die zich met het ijzer
verbindt: het ijzeroxide weegt evenveel als het ijzer plus die zuurstof.

\textbf{2.} Alles blijft in de pot: de verbruikte zuurstof zat al in de
pot, en het gevormde oxide blijft daar. De totale massa verandert niet.

\textbf{3.} \omterm{def:g8:balanced-equations:conservation}{Behoud van massa}: bij een reactie is de totale massa van de
\omterm{def:g7:chemical-reactions:reactant}{reactieproducten} gelijk aan de totale massa van de \omterm{def:g7:chemical-reactions:reactant}{beginstoffen}.

\textbf{4.} \ce{4Fe + 3O2 -> 2Fe2O3}.

\textbf{5.} \ce{3Fe + 2O2 -> Fe3O4}.

\textbf{6.} Links: 4 Fe, $3 \times 2 = 6$ O. Rechts: $2 \times 2 = 4$ Fe,
$2 \times 3 = 6$ O. Hij klopt.

\textbf{7.} 4 ijzeratomen op 3 zuurstofmoleculen, dus $400 \times
\frac{3}{4} = 300$ zuurstofmoleculen.

\textbf{8.} $0.28 \times \frac{3}{7} = \qty{0.12}{g}$ zuurstof.

\textbf{9.} $0.28 + 0.12 = \qty{0.40}{g}$ ijzeroxide.

\textbf{10.} Ja: er is \qty{0.12}{g} nodig en de pot bevatte ongeveer
\qty{0.5}{g}.

\textbf{11.} De verbruikte zuurstof heeft de \omterm{def:g4:air-a-mixture-of-gases:air}{lucht} in de pot verlaten
(hij zit nu in het vaste oxide), dus bevat de pot minder gas dan eerst
en stroomt er \omterm{def:g4:air-a-mixture-of-gases:air}{lucht} uit het lokaal naar binnen. De aflezing van de
balans stijgt dan, met de massa van de \omterm{def:g4:air-a-mixture-of-gases:air}{lucht} die binnenkomt.

\textbf{12.} Er is \qty{0.12}{g} zuurstof uit de \omterm{def:g4:air-a-mixture-of-gases:air}{lucht} in de pot
gehaald.
\end{solution}
